%=====================================================================
\section{Romanization}
%=====================================================================

The forced aligner used in the next stage was trained on romanized text, and it only recognizes the Latin letters a-z and the apostrophe. Each Bengali word is therefore converted to a Latin spelling by the romanizer (Table~\ref{tab:models}), and any remaining characters outside that set are removed. Words keep a link back to their original Bengali form, so that later stages can work with the native script again.

\begin{example}[: romanization]
The diagram shows each Bengali word above its romanized form. The two words outlined in red are English names written in Bengali script. \textbf{Neither romanization matches the English pronunciation}, and the doubled letters in \texttt{inglyaanndde} make the word look longer than it sounds.

\begin{center}
% Worked example, Stage 2: each Bengali word above its romanized form.
% Words are laid out in a grid of 6 columns; lossy romanizations are outlined in red.
\begin{tikzpicture}[
    cell/.style={rounded corners=2pt, minimum width=2.3cm, minimum height=1.35cm},
    ok/.style={draw=blue!55!black, text=blue!55!black},
    lossy/.style={draw=red!70!black, fill=red!6, text=red!70!black},
  ]
  \foreach [count=\i from 0] \bengali/\romanized/\note in \romanization {
    \pgfmathtruncatemacro\col{mod(\i,6)}
    \pgfmathtruncatemacro\row{div(\i,6)}
    \begin{scope}[shift={({\col*2.45},{-\row*1.55})}]
      \node[cell, \note] at (0,-0.2) {};
      \node at (0,0.08) {\bn \bengali};
      \draw[-{Stealth[length=1.5mm]}, \note] (0,-0.12) -- (0,-0.36);
      \node[font=\scriptsize\ttfamily, \note, draw=none, fill=none] at (0,-0.55) {\romanized};
    \end{scope}
  }
\end{tikzpicture}

\end{center}
\end{example}
